% Page budget: 1.55 pages | Owns: negation, OBC derivation, additional-state scope, implementation, and true-parameter condition.
% WRITING GUIDE
% Purpose: Explain negation, derive the final state-level orthogonal-by-construction map, and delimit its interpretability guarantee.
% Start here: Decisions D-185, D-190, D-192, D-194, D-198, and D-200 to D-203 in docs/decisions.md.
% Supporting material: Current state-level OBC derivations under scripts/gantry/orthogonal-by-construction/documentation and their thesis-readiness guidance.
% Mandatory papers: Read Györök's projection-regularization paper and orthogonal-by-construction paper in full before writing; keep their guarantees and assumptions separate.
% Research-plan anchor: Preserve Aspect 3's goal of protecting physical interpretability; state clearly how the final construction differs from the proposed regularisation.
% Current decisions: Reduced identifiable coordinates, data-derived reference tuples, one-step physical-state projection, per-epoch refresh, and a separate affine arm.
% Choices to justify: Protected parameter coordinates, reference-set construction, projected rows, rank tolerance, tangent versus affine basis, refresh frequency, and treatment of additional states.
% Implementation audit: Read model_augmentation/fit_systems/obc.py, scripts/gantry/gantry_dynamic/obc_gantry.py, obc_diagnostics.py, and the OBC attachment in model.py before fixing the derivation.
% Reference check: Verify the original projection method and separate its claims from the state-level, nonlinear-parameter, LPV, and additional-state extensions developed here.
% Cautions: docs/orthogonal-projection-plan.md describes an older soft-penalty route; one-step state orthogonality is not trajectory-output orthogonality or guaranteed true recovery.
% Planned figures: New geometric projection figure with a physical and additional-state partition inset; show the remaining uniqueness limitation as clearly as the protected direction.
% Section is finished when: The protected space, coordinate frame, gradient treatment, added-state scope, and true-recovery condition are explicit.
\section{Interpretability by Orthogonal Construction}

\subsection{Negation and the protected tangent space}
\begin{itemize}
\item Joint estimation is non-unique when the ANN can reproduce or cancel changes reachable by physical parameters \cite{gyorok2025l4dc,hoekstra2026lfraug}
\item Define interpretability as recovery of the ten identifiable physical combinations, not visual plausibility of individual raw parameters
\item Stack one-step sensitivities $J=\partial f_\theta/\partial\theta$ over a data-derived reference set in normalised state coordinates
\item Protect the rank-ten tangent space and compare the optional affine-offset column as a separate arm rather than silently changing the definition
\item Report tangent-space rank, condition number, and correlated sensitivity columns before interpreting parameter recovery
\item Figure (new): geometric decomposition $g=P_Jg+(I-P_J)g$ showing the baseline-aligned correction removed and the orthogonal correction retained
\item Use the same figure to show that local tangent separation does not imply global, trajectory-level, or latent-state uniqueness
\end{itemize}
\todo{Should the thesis call the affine arm a stricter non-negation definition or only a sensitivity analysis around the chosen tangent definition?}

\subsection{State-level orthogonal-by-construction map}
\begin{itemize}
\item Project the ANN write into the physical rows with $I-JJ^+$, eliminating the penalty weight used by soft orthogonality regularisation \cite{gyorok2025l4dc}
\item Freeze the reference basis within each epoch and refresh it as the physical parameters move, keeping the recurrent transition causal
\item Differentiate through the projection coefficient while treating the stacked basis as constant during each optimisation interval
\item Orthogonality applies to the physical one-step write; additional-state rows are orthogonal to the baseline by the baseline's zero extension
\item Additional states still propagate into later physical outputs, so one-step state orthogonality does not imply trajectory-output orthogonality
\item Separate the no-projection, ten-column tangent, and eleven-column affine arms under matched data, initialisation, and optimisation budgets
\item Figure change: add a small state-partition inset showing projection on physical rows, unrestricted additional-state rows, and their delayed influence through recurrent propagation rather than a direct output path
\end{itemize}
\todo{Is one-step state orthogonality the intended thesis contribution, or must the trajectory-level penalty also appear despite its unresolved weight selection?}

\subsection{Condition for true parameter recovery}
\begin{itemize}
\item OBC makes the local decomposition unique, but true recovery additionally requires the real missing dynamics $\Delta^*$ to satisfy $J^\top\Delta^*=0$
\item Derive the least-squares parameter bias $\hat\theta-\theta^*=J^+\Delta^*$ and define $\rho^*=\|JJ^+\Delta^*\|/\|\Delta^*\|$
\item Measure $\rho^*=0.205$ for the hidden absorber, so orthogonality fails on the current excitation and zero parameter error is not the valid target
\item The predicted combination-error floor is 4.96\%, dominated by a mass-imbalance direction sharing the absorber's equilibrium offset
\item Exact pointwise orthogonality is impossible for a nonzero force on this system; only excitation-averaged orthogonality can be designed
\end{itemize}
\todo{Should the proposed orthogonal multisine design be a future-work result, or must it be run before claiming a constructive remedy?}
